\section{Introduction}
\label{sec:introduction}

Large language models (LLMs) can now generate biomedical research abstracts
that are syntactically fluent, terminologically plausible, and difficult for
non-experts to distinguish from real PubMed entries
\cite{else2023chatgpt,liang2023gptdetect}. The downstream risks are concrete:
fabricated abstracts can pollute citation databases, support fraudulent
manuscripts, and amplify medical misinformation when they enter
patient-facing search results \cite{majovsky2023artificialintelligence}. A
biomedical-specific deepfake-text detector is therefore not a hypothetical
need but a near-term tooling gap for journals, preprint servers, and
clinical-information platforms.

A recurring promise in the AI-text-detection literature is that
\emph{adversarial training} — iteratively generating fakes, retraining the
detector against them, and repeating — yields a more robust binary
classifier \cite{zellers2019defending,mitchell2023detectgpt,hu2023radar}. We set out to test whether this promise survives a
controlled, fully reproducible biomedical pipeline running on commodity
hardware (a single Apple~M3~Max). Our pipeline pairs a
\textbf{BiomedBERT-Large} detector \cite{gu2021biomedbert} with two
generators (\textbf{BioMistral-7B} \cite{labrak2024biomistral} and
\textbf{SeqGAN} \cite{yu2017seqgan}) and an LLM rewriting
\textbf{agent} (\textbf{Qwen2.5-7B} \cite{qwen2024technical}) that targets
the top-\(25\%\) hardest fakes per round. Four conditions isolate the
contribution of each component (Section~\ref{sec:setup}).

\textbf{Research questions.} We ask:
\begin{itemize}
\item \textbf{RQ1.} Does 3-round adversarial retraining lower the evasion
  rate of a BiomedBERT detector against in-family fakes
  (BioMistral)?
\item \textbf{RQ2.} Does a Qwen2.5-7B rewriting agent meaningfully harden
  the fake distribution beyond the zero-shot generator?
\item \textbf{RQ3.} Does the resulting detector generalise to a
  cross-domain mixed-LLM benchmark (RAID \cite{dugan2024raid})?
\end{itemize}

\textbf{Contributions.} We treat the answers — all of them negative — as the
paper's contribution rather than an embarrassment to hide.
\begin{enumerate}[leftmargin=*]
\item \textbf{Generator-family dependence dominates.} Across one detector
  and one test set, evasion ranges from \(0.7\%\) (BioMistral fakes) to
  \(98.7\%\) (SeqGAN fakes) to \(98.3\%\) (RAID's mixed LLMs) — a
  \(\sim\!100\times\) spread that dwarfs every other effect we measured.
\item \textbf{Adversarial retraining produced no measurable improvement at
  \(3\)k scale.} Condition~B's per-round evasion stayed at
  \(0.993, 0.980, 0.987\) — within Wilson \(95\%\) CI noise on \(n=250\).
  Conditions~C and~D produced byte-identical test metrics because the
  best-validation-AUC checkpoint selector kept rolling back to the
  baseline.
\item \textbf{We retract a prior \(500\)-sample finding.} An earlier run
  reported a \(0.74 \to 0.54\) evasion curve in Condition~B; with the
  larger \(3\)k corpus and tighter CIs the curve disappears, and we now
  attribute it to small-sample noise.
\item \textbf{Cross-domain transfer is open.} A detector trained on
  PubMed~+~BioMistral fails on RAID (accuracy \(0.11\), evasion \(0.983\)),
  motivating future multi-family training rather than supporting any claim
  of universal AI-text detection.
\item \textbf{Reproducibility.} All code, data, metrics, plots, and a
  one-command rerun script are released; the entire study costs
  \(10\,\mathrm{h}\,01\,\mathrm{m}\) on a single laptop.
\end{enumerate}

\item \textbf{Adversarial-aware checkpoint selection.} We implement a
  drop-in replacement for best-val-AUC selection
  (\texttt{--checkpoint\_selector adversarial}) that scores all epoch
  checkpoints on the round's live fake pool and promotes the epoch with the
  \emph{lowest evasion rate}. This is designed to dissolve the C\,$=$\,D
  pathology when the val-AUC ceiling is already saturated.
\item \textbf{Multi-family training (Condition E).} We implement a
  BiomedBERT-Large detector trained jointly on BioMistral-7B, SeqGAN, and
  Llama-3.2-3B fakes. The multi-family corpus is assembled by
  \texttt{models/detector/train\_multi\_family.py} and is designed to
  compress the \(\sim\!140\times\) evasion spread identified in
  Contribution~1.
\item \textbf{Post-hoc calibration.} We implement temperature-scaling
  calibration~\cite{guo2017calibration}
  (\texttt{evaluation/calibration.py}) that learns an optimal decision
  threshold and temperature on a held-out calibration set. Applied to
  Condition~A's detector on the RAID probe, calibration recovers meaningful
  accuracy by correcting the threshold shift that produces the
  AUC--accuracy gap.

\textbf{Paper organisation.} Section~\ref{sec:related} surveys
machine-generated text detection, adversarial training for text classifiers,
and LLM-as-agent rewriting. Section~\ref{sec:method} describes the BioShield
pipeline. Section~\ref{sec:setup} details the four-condition ablation, data
splits, and metrics. Section~\ref{sec:results} reports per-condition results
and the RAID probe. Section~\ref{sec:discussion} owns the negative findings
and explains the C\,$=$\,D pathology. Section~\ref{sec:limitations}
catalogues limitations, and Section~\ref{sec:conclusion} closes.
